%% Dans ce fichier, les �ventuels accents sont cod�s
%% en ISO-8859-15 (aka latin9)
%% La ligne suivante donne trois exemples 
%% e accent aigu : �, e accent grave : �, e-dans-l'o : �
%%%%%ENTREES%%%%
\begin{tabarg}
\multicolumn{3}{l}{\titretab {Entr\'ees}}\\
\hline
$P_0$ & Vecteur, \texttt{float} &  Pression au niveau $Z_0$ (hPa)\\
$T_m$ & Vecteur, \texttt{float} & Temp�rature moyenne de la couche (Kelvin ou Celsius )\\
$Z_0$ & Vecteur, \texttt{float} & Altitude de la pression $P_0$ (m)\\
$Z$ & Vecteur, \texttt{float} & Altitude  (m) pour laquelle on calcule la pression\\
$R_a/g$ & Coeff., \texttt{float} & Rapport de la constante sp\'ecifique de
l'air sec sur l'acc\'el\'eration de la pesanteur\\
\end{tabarg}

\medskip
%%%%%SORTIES%%%%
\begin{tabarg}
\multicolumn{3}{l}{\titretab {Sorties}}\\
\hline
$P$ & Vecteur, \texttt{float} & Pression estim�e \`a l'altitude Z (hPa)\\
\end{tabarg}
